\section{Statement and conventions}
Let $\gamma_N$ be standard Gaussian measure on $\RR^N$ and define
\[
 T_\rho u(x)=\E u\bigl(\rho x+\sqrt{1-\rho^2}\,Y\bigr),
 \qquad Y\sim\gamma_N,\quad -1<\rho<1.
\]
An assignment is a measurable map $f=(f_1,\ldots,f_m)$ into
$\Delta_m=\{a\in[0,1]^m:\sum_i a_i=1\}$. Its noise stability is
$S_\rho(f)=\sum_i\ip{f_i}{T_\rho f_i}_{L^2(\gamma_N)}$.
This is a single-assignment functional; two independently chosen rules
are not permitted in its definition.

Let $P=I_m-\one\one^T/m$. Choose vectors $b_i^0\in\RR^N$ whose Gram
matrix is $\alpha^2P$, where
\begin{equation}\label{eq:scales}
\begin{aligned}
 n&=m-1,& L&=1+\log m,& \mu_m&=\E\max_{i\le m}Z_i,\\
 \alpha&=\frac{\mu_m}{m-1},& w&=\alpha^2(m-1),&
 \beta&=\sqrt{w/m}.
\end{aligned}
\end{equation}
Here the $Z_i$ are independent standard Gaussians. The regular assignment
$g$ chooses the largest $b_i^0\cdot x$. Ties are Gaussian null. Its
centroid matrix is $B_0=(b_i^0)_i$, and its cells have mass $1/m$.
The construction exists for every finite $N\ge m-1$ and extends
cylindrically when $N>m-1$.

\begin{theorem}\label{thm:main}
There is an absolute integer $m_0$ such that, for every $m\ge m_0$,
every finite $N\ge m-1$, every $\rho\in[-1/(m-1),0)$, and every
measurable assignment $f:\RR^N\to\Delta_m$,
\[
                         S_\rho(f)\ge S_\rho(g).
\]
Equality holds precisely for the indicator assignments of regular
Gaussian simplex partitions, up to orthogonal transformations,
relabeling, and Gaussian null sets. No mass constraint is imposed on $f$.
The threshold $m_0$ is independent of $\rho$ and of finite $N$.
\end{theorem}

This manuscript provides no numerical value of $m_0$. It does not claim
the positive-correlation conjecture, the remaining small values of $m$,
all negative Gaussian correlations, or a bilinear theorem. Section~\ref{sec:discrete}
gives the exact negative Plurality is Stablest consequence and the
UGC-conditional Frieze--Jerrum threshold. These are fixed-$m$ consequences,
with the endpoint formula credited to Garc\'ia-Soriano--Miyauchi \cite{GSM}.

The original negative Standard Simplex formulation has no mass constraint
\cite[Conjecture~1.4]{IM}. The older large-$m$ endpoint estimate in
\cite[Proposition~11.4]{KKMO}, recalled in \cite[equation~(20)]{IM}, is
asymptotic, with an additive error of order $\log\log(m)/m^2$.
Theorem~\ref{thm:main} instead asserts exact optimality for each fixed
sufficiently large $m$ on an entire interval. The bounded primary-source
check \cite{scope} is not an exhaustive priority or novelty determination.

For an assignment write
\begin{equation}\label{eq:assignment-data}
 p_i=\E f_i,\qquad B_i=\E[Xf_i],\qquad
 W=\norm{B}_{\Frob}^2,\qquad U=m^{-1}\one.
\end{equation}
Always $B^T\one=0$. For a centered score matrix $B$ define
\begin{align}
 F_p(B)&=\sup\left\{\E\sum_i f_i B_i\cdot X:
                  f:\RR^N\to\Delta_m,\ \E f=p\right\},\label{eq:dual-value}\\
 \calR(B)&=\frac{w+\norm{B}_{\Frob}^2}{2}-F_U(B).\label{eq:R-def}
\end{align}
Its value depends only on $BB^T$. Let
$\Gamma(C)=F_U(C^{1/2})$ for a centered positive semidefinite matrix
$C$; all square-root realizations are equivalent by Gaussian invariance.
The regular orbit distance is
$d(B)=\inf_{B_0}\norm{B-B_0}_{\Frob}$, over regular centroid matrices
in the same ambient space. Constants denoted $c,C$ may change from line
to line. A subscript records fixed parameter dependence. All estimates
called uniform below are independent of finite $N$ and of the correlation
in the stated interval.

\section{Analytic interfaces and dependency structure}\label{sec:inputs}
The following interfaces record the hypotheses used by the final assembly.
They are mathematical proof obligations, not consequences of an audit label.
Sections~\ref{sec:foundation-proof}, \ref{sec:local-proof}, and
\ref{sec:gap-proof} reconstruct their principal arguments; the accompanying
analytic notes listed in Appendix~\ref{app:sources} supply the longer
finite-fan, transport, and boundary calculations. Those notes form part of
the proof package and are required to assess the complete claim. The external
Gaussian comparison is separately identified below. The radial estimate in
Section~\ref{sec:radial} is not assumed as an interface.

\begin{prior}[Covariance value, normal strain, and offsets]
\label{prior:value}
The functional $\Gamma$ is concave and
continuous on centered positive semidefinite matrices. At a positive
definite matrix on $\one^\perp$, its derivative is $\calL/2$, where
$\calL$ is the mass-offset Laplacian. In particular
$\Gamma(P)=\mu_m$. It is $C^2$ locally in the positive cone. At
$C_0=\alpha^2P$, the normal-strain estimate implies
\begin{equation}\label{eq:reference-normal-strain}
 D^2\Gamma(C_0)[E,E]<0
 \quad\text{if some }(e_i-e_j)^TE(e_i-e_j)\ne0,
\end{equation}
for a centered symmetric $E$. Every regular facet has positive Gaussian
measure, so this strictness criterion applies to each pair.

If $\norm{B}_{\Frob}^2=w$, $\max_i\abs{B_i}\le A\alpha$, and $h_U$
is a mean-zero balanced optimal offset, then
\begin{equation}\label{eq:offset-input}
                         \norm{h_U}^2\le C_A L^2\calR(B).
\end{equation}
This is a value/offset statement at arbitrary rank, including coincident
sites by continuity. Positive scaling gives
\begin{equation}\label{eq:radial-identity}
 \calR(sB)=\tfrac12 w(s-1)^2+s\calR(B),\qquad
 h_U(sB)=s h_U(B)\quad(s>0,\ \norm{B}^2=w).
\end{equation}
In the mass box $p_i\in[1/(2m),3/(2m)]$, the optimum first-moment value
$V(p)=\sup_{\E f=p}\norm{M_f}^2$ satisfies
$V(p)\le w+CL^2\norm{p-U}^2$.
\end{prior}
The covariance and normal-strain input is supplied by the classical-input
Hodge argument \cite{hodge}, summarized in Section~\ref{sec:foundation-proof}.
The offset and mass-envelope argument additionally uses the stronger
cycle-refined Ehrhard inequality \cite{cycle,currentassembly}; this dependency
is not removed by the Hodge replacement.

\begin{proposition}[Endpoint and zero-remainder rigidity]\label{prop:endpoint}
For every centered $B$, $\calR(B)\ge0$, with equality precisely when
$BB^T=\alpha^2P$. Every balanced measurable assignment has $W\le w$;
equality forces the regular assignment. More generally, if an arbitrary
assignment's actual centroids have $\calR(B)=0$, it is regular.
\end{proposition}
\begin{proof}
We give the direct covariance proof \cite{endpoint}. For iid Gaussian
coordinates, the expected winning coordinate is $\mu_m$. Label symmetry
gives diagonal centroid entry $\mu_m/m$. The common coordinate mean is
independent of the winning label and has mean zero. The off-diagonal
entries are equal by symmetry, and centering makes the centroid matrix
exactly $B_0=\alpha P$ on $\one^\perp$. Its masses are $1/m$ and
$F_U(B_0)=\alpha\mu_m=w$.

The reference winner is unique off Gaussian-null hyperplanes. Hence the
score-value envelope is differentiable there and has derivative $B_0$.
Differentiating $C=BB^T$ at $B_0$, which is invertible on $\one^\perp$,
gives $\nabla\Gamma(C_0)=P/2$. Concavity therefore yields
\[
 \Gamma(C)\le w+\tfrac12(\tr C-w)
\]
for every centered positive semidefinite $C$, proving nonnegativity.
If $C\ne C_0$, put $E=C-C_0$. If every edge quadratic in
\eqref{eq:reference-normal-strain} vanished, writing $d_i=E_{ii}$
would give $E_{ij}=(d_i+d_j)/2$. The row sums then imply
$md_i+\tr E=0$ for each $i$, and summing forces $\tr E=0$ and $E=0$,
a contradiction. Thus the path $C_0+tE$ has strictly negative second
derivative at zero, and by local $C^2$ regularity on a short initial
interval. Its positive gap below the reference tangent is thereafter
nondecreasing by concavity. Continuity handles a singular endpoint at
$t=1$. This proves $\calR(B)>0$ unless $BB^T=C_0$.

For balanced $f$ with actual centroids $B$, feasibility gives
$W\le F_U(B)\le(w+W)/2$, hence $W\le w$. Equality implies
$BB^T=C_0$. For any assignment with this actual Gram matrix,
$W=w=\E\max_i B_i\cdot X$. Its nonnegative routing loss
$\max_i B_i\cdot X-\sum_i f_i B_i\cdot X$ has expectation zero.
The regular rows have Gaussian-null ties, so $f$ is their regular
winning assignment. This final implication needs no mass assumption.
\end{proof}

\begin{prior}[Competitive localization]
\label{prior:localization}
For all sufficiently large $m$ and $0<r\le1/(m-1)$, an assignment with
$S_{-r}(f)\le S_{-r}(g)$ has $p_i\in[1/(2m),3/(2m)]$. Put
\[
 \delta=\norm{p-U}^2,\quad \Delta=w-W,\quad D=\delta+r\Delta.
\]
Then
\begin{equation}\label{eq:localization}
              0\le\delta\le2D,\qquad 0\le D\le Cr^2/m.
\end{equation}
\end{prior}
The proof is reconstructed in Section~\ref{sec:foundation-proof} from
\cite{localization,currentassembly}; it applies to arbitrary competitive assignments. For every assignment,
conditional Cauchy gives
$(B_i\cdot v)^2\le p_i\E[f_i(X)(v\cdot X)^2]$. Summation shows
\begin{equation}\label{eq:actual-operator}
                              B^TB\preceq p_{\max}I.
\end{equation}
Conditional Jensen in $\E e^{t v\cdot X}$ similarly gives
\begin{equation}\label{eq:entropy-centroid}
                 \abs{B_i}\le p_i\sqrt{2\log(1/p_i)}.
\end{equation}
Neither estimate requires stationarity or purity.

\begin{prior}[Global one-row covariance coercivity]\label{prior:one-row}
For every fixed $A$, every centered score matrix with
$\norm{B}^2=w$ and $\max_i\abs{B_i}\le A\alpha$ satisfies
\begin{equation}\label{eq:one-row}
 \max_i\left|\frac{\abs{B_i}^2}{\beta^2}-1\right|^2
                         \le C_A\frac{m^2}{L}\calR(B).
\end{equation}
No rank, assignment-realizability, or minimization hypothesis is imposed.
\end{prior}
This is the global single-row estimate of \cite[Theorem~1]{rowcontrol}.
It will be applied directly to a normalized actual centroid matrix;
no actual-minimizer row theorem is assumed.

\begin{prior}[Global covariance remainder and local full-noise comparison]
\label{prior:local-global}
For centered $B$, put $C=BB^T/\alpha^2$ and $e=\norm{C-P}_{\Frob}$.
The initial-segment covariance bound is
\begin{equation}\label{eq:global-covariance-floor}
 \calR(B)\ge \frac{c\alpha^2e^2}{m}
                   \min\left\{1,\frac{c}{Le}\right\}.
\end{equation}
It holds at arbitrary endpoint rank and does not require
$\norm{B}^2=w$; the value at $e=0$ is interpreted by continuity.

For every fixed $A$, sufficiently large $m$, and every measurable
assignment whose \emph{actual} centroids obey $d(B)\le A/m$ and whose
masses lie in $[1/(2m),3/(2m)]$,
\begin{equation}\label{eq:local-noise}
 S_{-r}(f)-S_{-r}(g)\ge \tfrac12\norm{p-U}^2+r\calR(B),
                   \qquad0<r\le1/(m-1).
\end{equation}
Competitiveness in this statement forces a regular assignment.
\end{prior}
Section~\ref{sec:gap-proof} proves the global ray estimate from the
quantitative Hodge input and the facet bootstrap \cite{currentgap}.
Section~\ref{sec:local-proof} reconstructs the local theorem
\cite{currentlocal}. Its complete cap and affine-boundary calculations are
in the accompanying source notes. It concerns arbitrary measurable
assignments, including fractional rules.

\begin{prior}[Regular normalization and margin]
\label{prior:regular}
Let $\Pi_k$ denote the orthogonal Hermite projection of degree $k$ and
$A_k(f)=\sum_i\norm{\Pi_k f_i}_2^2$. Then
\begin{equation}\label{eq:A2-reference}
 A_2(g)=\frac{m\{\E[(\max_i Z_i)^2]-1\}^2}{2(m-1)(m-2)},
 \qquad \alpha^2\asymp L/m^2.
\end{equation}
For a regular score matrix $B_0$, let $G(X)$ be its largest score minus
its second largest score. Then
\begin{equation}\label{eq:regular-margin}
                        \Prob\{G\le t\}\le Cmt\quad(t\ge0).
\end{equation}
\end{prior}
The exact second-chaos identity is in \cite{localization}; the margin
and its fractional routing consequence are in \cite{globaltail}.
We also use the standard dimension-free Gaussian hypercontractivity
$\norm{T_tu}_2\le\norm{u}_{1+t^2}$ \cite{Gross}, and the one-dimensional variance
bound $\Var_\nu x\le1/c$ when the density of $\nu$ has logarithmic
second derivative at most $-c$ \cite{BrascampLieb}. The transitive proofs
of the offset and mass-window inputs also use Gaussian isoperimetry
\cite{Borell} and the convex Ehrhard inequality \cite{Latala}; their
precise uses are recorded in the companion section-level source map.

\begin{prior}[Equal-threshold Gaussian maximum comparison]
\label{prior:max}
If $K$ is an $m\times m$ correlation matrix and
$K-\one\one^T/m\succeq0$, then for $Y\sim N(0,K)$ and independent
standard Gaussians $Z_i$,
\[
                 \Prob\{\max_iY_i\le t\}\ge\Phi(t)^m
                 \qquad(t\in\RR).
\]
The statement includes singular $K$.
\end{prior}
This is Theorem~2.1 of the pinned primary source \cite{Mulgund}. Only
equal thresholds are used. We do not infer a general rectangular
comparison or a comparison for unequal marginal variances.
